Medical image analysis models need to adapt as medical centers, disease categories, organs, and imaging modalities change, while retaining their performance on earlier tasks. The information available for training is often limited: detailed segmentation annotations are costly, raw images from different centers cannot readily be pooled, and historical samples may no longer be accessible. In registration, changes in organs and modalities also alter registration objectives and deformation distributions, making adaptation and retention difficult to balance with only a few historical image pairs. This thesis addresses these problems by designing a continual learning evaluation platform and developing methods for weakly supervised segmentation, replay-free federated classification, and registration with bounded replay.

To compare continual learning under different training conditions, this thesis designs the MedCL evaluation platform. MedCL covers segmentation, classification, and registration, and separately records analysis tasks, incremental settings, and available training information. It associates locally trained stage models or their predictions with fixed test data. Stage--task performance matrices distinguish final performance, historical retention, current-task learning, and forward generalization where applicable. Model parameters and historical-information usage records provide complementary measures of resource cost. Task adapters preserve the meaning of each output and the direction and units of its metrics. This provides a consistent way to organize and analyze subsequent experiments without directly combining scores from different tasks.

For segmentation, this thesis constructs a continual learning benchmark covering cross-center domain shifts, incremental cardiac structures, and cross-organ tasks. Representative methods are compared using shared data splits, network architectures, and training settings. The results show that reducing forgetting does not necessarily improve new-task learning or generalization to future domains; performance also depends on parameter growth, replay capacity, and task order. To address insufficient annotations and forgetting together, the thesis then proposes ZSDERpp. Global consistency and spatial priors provide training constraints for unlabeled regions, while bounded replay of image--scribble pairs and reference-feature constraints help retain old-task performance. MiB is incorporated to handle changing background semantics in the class-incremental setting. Weakly supervised experiments reuse the benchmark's patient splits, task orders, and dense test references. Experiments on six-center domain-incremental and three-stage class-incremental segmentation show improvements in segmentation performance and historical retention. Future-domain generalization, evaluated in the domain-incremental experiment, does not improve correspondingly.

For federated continual classification without historical raw-sample replay, this thesis proposes FedSubMerge, which fuses principal gradient subspaces across clients. Clients represent historical gradient directions with low-rank subspaces, and the server combines their protection information. Subsequent local updates are projected onto the orthogonal complement of the protection space to reduce interference with old tasks. An adaptive version, FedSubMerge-AD, selects related clients layer by layer, reducing constraints imposed by unrelated historical directions. Experiments use sequentially arriving class batches and task-specific classification heads with known test-time task identity. Evaluation covers distribution skew, label-quantity skew, and source-related feature skew. Both fusion strategies improve final classification performance without replaying historical raw samples, while their relative advantages vary with client distributions.

For registration as organs and modalities change, this thesis proposes SAMCL, combining bounded experience replay, first-order meta-continual learning, and sharpness-aware optimization. A small set of historical fixed--moving image pairs supplies old-task training signals. Meta-continual learning coordinates updates from current and historical samples, while sharpness-aware optimization considers loss variation in the parameter neighborhood. Experiments use a sequence of brain, abdominal, lung, and cross-modal abdominal registration tasks. Following the platform's stage-based evaluation, structural Dice and target registration error (TRE) are analyzed separately to assess registration accuracy, historical retention, and current-task learning. Further experiments examine replay capacity and within-domain and cross-task generalization. SAMCL reduces historical degradation compared with direct sequential training. Compared with meta-experience replay, it performs better on structural-overlap measures, while landmark-error measures show a different trade-off.
